% Canonical 20-step full-network SGD loss, from scripts/teaching_trace.py.
% Checkpoints reproduce teaching_reference.py's initial and final states.
\tagstructbegin{tag=Figure,alttext={三维教学模型零到二十步联合梯度下降的平均交叉熵下降曲线。初始为一点七九八四三零，第一步一点六一一三七七，第九步约零点八零三九一八，第二十步约零点四八三一七三。}}
\begin{center}
\begin{tikzpicture}
\begin{axis}[
width=116mm,height=58mm,
xlabel={全参数 SGD 步数},ylabel={三位置平均交叉熵},
xmin=0,xmax=20,ymin=.35,ymax=1.9,
xtick={0,1,5,10,15,20},ytick={.5,1,1.5},
tick label style={font=\footnotesize},label style={font=\small},
grid=major,grid style={linegray},
]
\addplot[blue,very thick,mark=*,mark size=1.5pt] coordinates {
(0,1.798430) (1,1.611377) (2,1.448561) (3,1.308360)
(4,1.187885) (5,1.084851) (6,.997085) (7,.922397)
(8,.858656) (9,.803918) (10,.756501) (11,.715015)
(12,.678708) (13,.646285) (14,.617030) (15,.590375)
(16,.565869) (17,.543156) (18,.521948) (19,.502016)
(20,.483173)};
\end{axis}
\end{tikzpicture}
\captionof{figure}{同一三维主例连续 20 步：所有九组可训练参数同步更新，非仅训练输出头}
\label{fig:canonical-20-step-loss}
\end{center}
\tagstructend
